\chapter{Reference: lists and selection}
\label{ch:refman-list}

\begin{aside}
A scrolling list of items with a highlighted selection.
Build it from rows and a model index; \maya{} ships nothing
fancier than that because nothing fancier is needed.
\end{aside}

\section{Working example}

\begin{lstlisting}[language=C++]
#include <maya/maya.hpp>
#include <vector>

using namespace maya;
using namespace maya::dsl;

struct ListDemo {
    struct Model {
        std::vector<std::string> items{
            "alpha", "beta", "gamma", "delta",
            "epsilon", "zeta", "eta", "theta"
        };
        int selected = 0;
    };
    struct Up {};
    struct Down {};
    struct Choose {};
    struct Quit {};
    using Msg = std::variant<Up, Down, Choose, Quit>;

    static Model init() { return {}; }

    static auto update(Model m, Msg msg)
        -> std::pair<Model, Cmd<Msg>>
    {
        return std::visit(overload{
            [&](Up) {
                m.selected = std::max(0, m.selected - 1);
                return std::pair{m, Cmd<Msg>{}};
            },
            [&](Down) {
                int n = static_cast<int>(m.items.size());
                m.selected = std::min(n - 1, m.selected + 1);
                return std::pair{m, Cmd<Msg>{}};
            },
            [&](Choose) { return std::pair{m, Cmd<Msg>{}}; },
            [](Quit) {
                return std::pair{Model{}, Cmd<Msg>::quit()};
            },
        }, msg);
    }

    static Element view(const Model& m) {
        std::vector<Element> rows;
        for (size_t i = 0; i < m.items.size(); ++i) {
            auto label = text(m.items[i]);
            if (static_cast<int>(i) == m.selected) {
                rows.push_back(
                    h(t<"> ">, label | Bold) | Bg<40, 80, 140>);
            } else {
                rows.push_back(h(t<"  ">, label));
            }
        }
        return v(
            t<"pick one"> | Bold,
            blank_,
            v(rows),
            blank_,
            t<"up/down move, enter choose, q quit"> | Dim
        ) | pad<1> | border_<Round>;
    }

    static auto subscribe(const Model&) -> Sub<Msg> {
        return key_map<Msg>({
            {'q', Quit{}},
            {SpecialKey::Up,    Up{}},
            {SpecialKey::Down,  Down{}},
            {SpecialKey::Enter, Choose{}},
        });
    }
};

int main() { run<ListDemo>({.title = "list"}); }
\end{lstlisting}

\section{Notes}

\begin{itemize}[leftmargin=*]
  \item For long lists, clip a window of items around the
    selection; rendering off-screen rows costs.
  \item Use \code{Bg<...>} on the selected row; readers
    parse highlight faster than markers.
  \item Pair with \code{home}/\code{end} for jump-to-edges.
\end{itemize}

\section{Recap}

\begin{recap}
\begin{itemize}[leftmargin=*]
  \item Build rows; mark one selected.
  \item Up / Down adjust a single index.
  \item Clip to a window for large lists.
\end{itemize}
\end{recap}

\section*{Workshop}

\begin{enumerate}[leftmargin=*]
  \item Add a filter input that prunes the visible items.
  \item Show the index and total in the footer.
  \item Implement multi-select with \code{space} to toggle.
\end{enumerate}
